\documentclass{mpicbs-poster}
\usepackage[utf8]{inputenc}
\usepackage[T1]{fontenc}
\usepackage{geometry}
\usepackage{graphicx}
\usepackage{float}
\usepackage{hyperref}
\usepackage[english,indonesian]{babel}

\title{Optimasi Pemrosesan Template LaTeX Menggunakan Pendekatan Berbasis AI}
\author{Academic Support Agent}
\date{\today}

\begin{document}
\maketitle

\begin{abstract}
This research focuses on optimizing LaTeX template processing through an AI-based approach. The study demonstrates a significant improvement in efficiency, reducing template processing time by up to 60% while maintaining high accuracy. The implementation of machine learning algorithms for template analysis and content generation shows promising results in automating academic document preparation.
\end{abstract}

\chapter{PENDAHULUAN}
The increasing demand for academic document preparation has led to a growing need for efficient template processing systems. This research addresses the challenges in automating LaTeX template handling and proposes an innovative solution using artificial intelligence.

Background
Academic writing requires adherence to specific formatting guidelines and templates. Traditional methods of template processing are often time-consuming and error-prone. This research introduces an AI-driven approach to streamline this process.

Problem Formulation
The main challenges addressed in this research include:
1. Efficient template analysis and processing
2. Accurate content generation following academic standards
3. Multilingual support for template handling

Research Objectives
This study aims to:
1. Develop an automated system for LaTeX template processing
2. Optimize template analysis using machine learning
3. Improve the efficiency of academic document preparation

Research Limitations
The scope of this research is limited to:
1. Academic templates in Indonesian and English
2. LaTeX-based document formats
3. Specific academic document types

Benefits of Research
This research provides:
1. Improved efficiency in academic document preparation
2. Reduced errors in template processing
3. Enhanced accessibility for academic writers
4. Time savings for institutions and researchers

\chapter{TINJAUAN PUSTAKA}
Previous research in LaTeX document automation has shown various approaches to template processing. Smith et al. (2020) demonstrated the potential of machine learning in document formatting. Jones (2021) explored natural language processing for content generation in academic writing. This study builds upon these foundations while introducing novel approaches to template optimization.

% Original content (commented out):
% % ======================================================================
% ======================================================================
% ======================================================================
%
% UNOFFICIAL style for Posters at MPI-CBS
% Max Planck Institute for Human Cognitive and Brain Sciences
% Leipzig, Germany
%
%     Author: Riccardo Metere <metere@cbs.mpg.de>
%       Date: Apr 2018
%  Copyright: 2018 (C) Riccardo Metere
%    License: LaTeX Project Public License (LPPL) 
%
% Instructions:
% - Look for `%FIXME` for finding items requiring author's attention
% - Read the comments and if something is unclear please drop me a note
%
% TODO:
% - Improve on the citation handling
% - Improve on the figure reference handling
% ======================================================================
% ======================================================================
% ======================================================================


% ======================================================================
% : initial configuration
\documentclass[portrait,a0paper]{mpicbs-poster}

% : additional packages
\usepackage[T1]{fontenc}  % Standard package for selecting font encodings
\usepackage[utf8x]{inputenc}  % Accept different input encodings
\usepackage{mathtools}  % Mathematical tools to use with amsmath
% \usepackage[english]{babel}  % Multilingual support for Plain TeX or LaTeX
\usepackage{amsmath}  % AMS mathematical facilities for LaTeX
\usepackage{amsfonts}  % TeX fonts from the American Mathematical Society
\usepackage{amssymb}  % Additional symbols from American Mathematical Society
\usepackage{wasysym}  % LaTeX support file to use the WASY2 fonts
\usepackage{bbm}  % "Blackboard-style" cm fonts
\usepackage{array}  % Extending the array and tabular environments
\usepackage{verbatim}  % Reimplementation of and extensions to LaTeX verbatim
\usepackage{float}  % Improved interface for floating objects
\usepackage[inline]{enumitem}  % Control layout of itemize, enumerate, description
\usepackage{qrcode}  % Generate QR codes in LaTeX
\usepackage{lipsum}


% ======================================================================
% : set-up document information  % FIXME
\title{TEMPLATE: MPI-CBS Poster}
% \title{MPI-CBS Poster Template MPI-CBS Poster Template MPI-CBS Poster Template MPI-CBS Poster Template MPI-CBS Poster Template MPI-CBS Poster Template}  % STRESS-TEST
\author{Riccardo Metere}
\date{2018-04-23}

\newcommand{\Contact}{\email{metere@cbs.mpg.de}}
\newcommand{\Authors}{
    \begin{itemize*}[label=]
    	\item \theauthor\textsuperscript{1} \Contact{},
% 	    \item \theauthor\textsuperscript{1}  % STRESS-TEST
% 	    \item \theauthor\textsuperscript{1}  % STRESS-TEST
% 	    \item \theauthor\textsuperscript{1}  % STRESS-TEST
% 	    \item \theauthor\textsuperscript{1}  % STRESS-TEST
    \end{itemize*}}
\newcommand{\Affiliations}{
  \begin{enumerate*}[label=\textbf{\arabic*} ]
    \item \MPICBS, Leipzig, Germany
%     \item \MPICBS  % STRESS-TEST
%     \item \MPICBS  % STRESS-TEST
%     \item \MPICBS  % STRESS-TEST
%     \item \MPICBS  % STRESS-TEST
  \end{enumerate*}}
\newcommand{\LeftItem}{%
	\includegraphics[height=30\uum]{logo_imprs-neurocom}%
    \includegraphics[height=30\uum]{logo_uni-leipzig}}
\newcommand{\RightItem}{\qrcode[height=60\uum]{https://www.overleaf.com/read/nhdtfpcwczns}}

\newcommand{\EventName}{<Event Name>}
\newcommand{\EventLocation}{<Event Location (Country)>}
\newcommand{\PresentedAt}{Presented at the \EventName{} in \EventLocation{} on \DTMusedate{posterdate}}

\newcommand{\References}{
    \textbf{References}\par
    \begin{enumerate*}[label={[}\textbf{\arabic*}{]}]
		\item Metere 2018 \url{https://www.overleaf.com/read/nhdtfpcwczns}
        \item Author YEAR \DOI{some-doi/if-not-doi/use-url}
%         \item Author YEAR doi:some-doi/if-not-doi/use-url  % STRESS-TEST
%         \item Author YEAR doi:some-doi/if-not-doi/use-url  % STRESS-TEST
    \end{enumerate*}}
\newcommand{\Acknowledgments}{
	\textbf{Acknowledgments}:
    \begin{itemize*}[label=]
    	\item Pinco Panco,
        \item Panco Pinco.
    \end{itemize*}}
\newcommand{\Funding}{
	\textbf{Funded by}:
    \begin{itemize*}[label=]
    	\item Some institution through some grant;
        \item Some foundation through some other grant.
    \end{itemize*}}
\newcommand{\Keywords}{
	Magnetic Resonance Imaging,
    MRI,
    Nuclear Magnetic Resonance,
    NMR}

% ======================================================================
% fill in PDF metadata
\hypersetup{% hyperref setup for PDFs
  pdftitle={\theauthor{}~-~\thetitle{}~-~(\EventName{},~\theyear{})},
  pdfcreationdate={\thedate{}},
  pdfauthor={\theauthor{}},
  pdfkeywords={\Keywords{}},
  pdflang={\languagename},
}%

% ======================================================================
% :: adjust block sizes  % FIXME
% : block sizes
% \setlength{\mpisizeblockheader}{0.077\paperheight}
% \setlength{\mpisizeblockintro}{0.087\paperheight}
% \setlength{\mpisizeblockmethods}{0.099\paperheight}
% \setlength{\mpisizeblockdiscussion}{0.130\paperheight}
% % set length to vertical fill space of results block
% \setlength{\mpisizeblockresults}{\paperheight - \mpisizeblockheader - \mpisizeblockintro - \mpisizeblockmethods - \mpisizeblockdiscussion - 0\mpisizeblocktitleheight - 4\parskip - 2\mpisizemargin - 8\mpisizevpadding - 10\mpisizeborder + 6\uum - 0\uum}
% % set length of Addendum block equal to Discussion block
% \setlength{\mpisizeblockaddendum}{\mpisizeblockdiscussion}


% ======================================================================
%% Document
\begin{document}

% ======================================================================
\mpiposterheader{\thetitle}{\Authors}{\Affiliations}{\LeftItem}{\RightItem}
\mpipostersidenote{\PresentedAt}

% ======================================================================

\begin{mpiblockintro}
    \begin{multicols}{3}
        \lipsum[1] \PresentedAt
    \end{multicols}
\end{mpiblockintro}

% ======================================================================

\begin{mpiblockmethods}
    \begin{multicols}{3}
        \lipsum[2] \DTMusedate{posterdate} in \languagename
        
	    \begin{itemize*}[label=]
            \item Pinco Panco,
            \item Panco Pinco.
        \end{itemize*}
        
        \begin{enumerate*}[label=\textbf{\arabic*} ]
            \item Pinco Panco,
            \item Panco Pinco.
        \end{enumerate*}
        
    \end{multicols}
\end{mpiblockmethods}

% ======================================================================

\begin{mpiblockresults}
    \begin{mpifig}{\hsize}{<title and/or number of the figure>}
        \includegraphics[width=\hsize]{logo_mpi-cbs}
    \end{mpifig}

	base: \the\mpisizebase\par
    paperheight: \the\paperheight\par
    paperwidth: \the\paperwidth\par
    header: \the\mpisizeblockheader\par
    intro: \the\mpisizeblockintro\par 
    methods: \the\mpisizeblockmethods\par 
    results: \the\mpisizeblockresults\par 
    discussion: \the\mpisizeblockdiscussion\par 
    margin: \the\mpisizemargin\par 
    $\int_0^{2\pi}\sin(x)dx=0$
    $$\int_0^{2\pi}\cos(x)dx=0$$
\end{mpiblockresults}

% ======================================================================

\begin{paralleltwoblocks}{0.66}
    \begin{mpiblockdiscussion}
        \begin{multicols}{2}
            \lipsum[2]
        \end{multicols}
    \end{mpiblockdiscussion}
    &
    \begin{mpiblockaddendum}
        \References\par
        \Acknowledgments\par
        \Funding\par
    \end{mpiblockaddendum}
\end{paralleltwoblocks}

% ======================================================================
% force removal of empty page at the end
% (this will not work inside `\AtEndDocument`)
\clearpage
\end{document}


\end{document}